\documentclass[12pt,a4paper]{article}

% ─────────────────────────────────────────
%  Packages
% ─────────────────────────────────────────
\usepackage[utf8]{inputenc}
\usepackage[T1]{fontenc}
\usepackage[english]{babel}
\usepackage{amsmath,amssymb,amsthm}
\usepackage{geometry}
\usepackage{graphicx}
\usepackage{booktabs}
\usepackage{array}
\usepackage{listings}
\usepackage{xcolor}
\usepackage{hyperref}
\usepackage{algorithm}
\usepackage{algpseudocode}
\usepackage{caption}
\usepackage{setspace}
\usepackage{titlesec}
\usepackage{fancyhdr}
\usepackage{enumitem}
\usepackage{tcolorbox}
\usepackage{float}

% ─────────────────────────────────────────
%  Page layout
% ─────────────────────────────────────────
\geometry{left=3cm, right=2.5cm, top=2.5cm, bottom=2.5cm}
\setlength{\parindent}{0pt}
\setlength{\parskip}{6pt}
\onehalfspacing

% ─────────────────────────────────────────
%  Header / Footer
% ─────────────────────────────────────────
\pagestyle{fancy}
\fancyhf{}
\lhead{\small Laporan Tugas}
\rhead{\small Perancangan dan Analisis Algoritma}
\cfoot{\thepage}
\renewcommand{\headrulewidth}{0.4pt}

% ─────────────────────────────────────────
%  Theorem environments
% ─────────────────────────────────────────
\theoremstyle{definition}
\newtheorem{definition}{Definition}[section]
\newtheorem{theorem}{Theorem}[section]
\newtheorem{observation}{Observation}[section]

% ─────────────────────────────────────────
%  Code listing style
% ─────────────────────────────────────────
\lstset{
  basicstyle=\ttfamily\footnotesize,
  keywordstyle=\color{blue}\bfseries,
  commentstyle=\color{gray}\itshape,
  stringstyle=\color{red!70!black},
  numberstyle=\tiny\color{gray},
  numbers=left,
  numbersep=6pt,
  frame=single,
  breaklines=true,
  showstringspaces=false,
  tabsize=4,
  captionpos=b
}

% ─────────────────────────────────────────
%  Coloured boxes for key insights
% ─────────────────────────────────────────
\tcbuselibrary{skins,breakable}
\newtcolorbox{keyinsight}[1][Key Insight]{
  title=#1,
  colback=blue!5!white,
  colframe=blue!50!black,
  fonttitle=\bfseries,
  breakable
}
\newtcolorbox{notebox}[1][Note]{
  title=#1,
  colback=yellow!10!white,
  colframe=orange!70!black,
  fonttitle=\bfseries,
  breakable
}
\newtcolorbox{connectionbox}[1][Koneksi CLRS–OAE]{
  title=#1,
  colback=green!5!white,
  colframe=green!50!black,
  fonttitle=\bfseries,
  breakable
}

% ─────────────────────────────────────────
%  Hyperref setup
% ─────────────────────────────────────────
\hypersetup{
  colorlinks=true,
  linkcolor=blue!60!black,
  urlcolor=blue!70!black,
  citecolor=green!50!black
}

% ═══════════════════════════════════════════
%  DOCUMENT BEGIN
% ═══════════════════════════════════════════
\begin{document}

% ─────────────────────────────────────────
%  Cover page
% ─────────────────────────────────────────
\begin{titlepage}
  \centering
  \vspace*{2cm}
  {\LARGE\bfseries Resume Materi Perkuliahan\\[0.4em]
   Perancangan dan Analisis Algoritma - A\par}
  \vspace{2cm}
  {\large Gilbran Mahdavikia Raja\\[0.3em]
   5025241134\\[20em]
   Teknik Informatika\\[0.3em]
   Institut Teknologi Sepuluh Nopember\par}
  \vspace{1cm}
  {\large 2026\par}
  \vfill
\end{titlepage}

% ─────────────────────────────────────────
%  Table of contents
% ─────────────────────────────────────────
\tableofcontents
\newpage

% ═══════════════════════════════════════════
\section{Pendahuluan}
% ═══════════════════════════════════════════

dalam perkuliahan Perancangan dan Analisis Algoritma (PAA) memperkenalkan dua strategi
kunci dalam merancang algoritma: \emph{Incremental} serta \emph{Divide and Conquer}.
Laporan ini membahas secara rinci implementasi metode Divide and Conquer untuk
memecahkan dua persoalan latihan yang melibatkan relasi rekurensi, yaitu:

\begin{enumerate}
  \item \textbf{SPOJ 6172 -- OAE (Only Altogether Exciting):} Persoalan kalkulasi
        banyaknya kombinasi bilangan bulat valid sepanjang $n$ digit yang bebas 
        dari kemunculan digit 0. Persoalan ini dituntaskan menggunakan prinsip 
        rekurensi dan pencarian solusi bentuk tertutup.
  \item \textbf{EOlymp -- New Year Tree:} Soal tentang penghitungan banyaknya 
        varian pewarnaan sebuah pohon Natal yang tersusun atas $n$ tingkat 
        dengan ketersediaan $k$ warna, yang dapat diselesaikan memanfaatkan 
        model rekurensi pohon.
\end{enumerate}

Pemilihan kedua persoalan tersebut didasari oleh alasan bahwa skema penyelesaiannya
jauh lebih natural jikalau dikerjakan dengan pendekatan rekurensi ketimbang 
incremental. Hal ini sangat sejalan dengan konsep yang dibahas di dalam buku 
\emph{Introduction to Algorithms} (CLRS), khususnya pada Bab~4 terkait metode 
Divide-and-Conquer.

% ═══════════════════════════════════════════
\section{Landasan Teori: Divide and Conquer dan Rekurensi}
% ═══════════════════════════════════════════

\subsection{Paradigma Divide and Conquer}

Merujuk pada Cormen et al.~\cite{clrs}, paradigma algoritma Divide and Conquer 
mencakup tiga prosedur esensial yang dieksekusi pada setiap kedalaman rekursi:

\begin{enumerate}
  \item \textbf{Divide:} Memecah permasalahan awal ke dalam sekumpulan submasalah 
        yang mana masing-masingnya merupakan versi yang lebih sempit dari persoalan tersebut.
  \item \textbf{Conquer:} Menuntaskan berbagai submasalah tersebut menggunakan langkah 
        rekursif. Ketika ukuran submasalah sudah sedemikian kecil, cari penyelesaiannya 
        secara langsung (dikenal sebagai \emph{base case}).
  \item \textbf{Combine:} Merangkai beragam solusi yang didapat dari setiap submasalah 
        guna membentuk konklusi akhir dari masalah semula.
\end{enumerate}

\subsection{Rekurensi (\textit{Recurrences})}

\begin{definition}
  Dalam terminologi matematika, rekurensi (\emph{recurrence}) merujuk pada persamaan 
  ataupun pertidaksamaan yang menyajikan spesifikasi suatu fungsi bersandarkan pada 
  penggunaan nilai fungsi tersebut di ukuran masukan yang lebih kecil~\cite{clrs}.
\end{definition}

Seperti yang telah diuraikan dalam referensi CLRS halaman 65--67 (Bab~4), rekurensi dan 
pendekatan divide-and-conquer bagaikan dua entitas yang saling melengkapi. Konsep 
rekurensi menyediakan medium yang sangat intuitif guna menjabarkan estimasi kompleksitas 
waktu eksekusi dari algoritma berjenis divide-and-conquer. Sebagai contoh, fungsi 
rekurensi klasik dari perancangan \textsc{Merge-Sort} adalah:

\begin{equation}
  T(n) =
  \begin{cases}
    \Theta(1)           & \text{jika } n = 1,\\
    2T(n/2) + \Theta(n) & \text{jika } n > 1.
  \end{cases}
\end{equation}

Ada tiga pendekatan primer untuk menyelesaikan rumusan rekurensi~\cite{clrs}:

\begin{itemize}
  \item \textbf{Metode Substitusi (\textit{Substitution Method}):} Menghipotesiskan
        sebuah batasan atas atau bawah, yang selanjutnya dibuktikan dengan 
        logika induksi matematis.
  \item \textbf{Metode Pohon Rekursi (\textit{Recursion-Tree Method}):}
        Mentransformasikan permasalahan rekurensi menjadi graf pohon, dengan 
        setiap node-nya mencerminkan representasi biaya komputasional pada tiap 
        tingkatan rekursif.
  \item \textbf{Metode Master (\textit{Master Method}):} Pendekatan praktis guna 
        mencari asumsi asimptotik (\emph{bounds}) dari relasi rekurensi dengan format 
        $T(n) = aT(n/b) + f(n)$.
\end{itemize}

Di dalam laporan ini, pengerjaan yang dipilih ialah metode yang menyertakan 
pencarian pola (\emph{pattern matching}) serta pengerjaan secara substitusi berulang 
untuk mendapatkan rumus \emph{closed-form solution} (solusi format tertutup).

% ═══════════════════════════════════════════
\section{Soal 1: SPOJ 6172 -- OAE (Only Altogether Exciting)}
% ═══════════════════════════════════════════

\subsection{Deskripsi Masalah}

\begin{tcolorbox}[title=SPOJ 6172 -- OAE, colframe=black!60, colback=gray!5]
  \textbf{URL:} \url{https://www.spoj.com/problems/OAE/}\\[4pt]
  Diketahui sebuah nilai integer bernilai logis positif berupa angka $N$. Cari tahu 
  total banyaknya bilangan real positif yang disusun oleh persis $N$ buah digit
  dan dipastikan \textbf{secara absolut terbebas dari angka 0}.
  Lakukan pencetakan terhadap jawaban modulus $314159$.\\[4pt]
  \begin{itemize}
    \item \textbf{Masukan:} Banyaknya ujicoba (test case) sejumlah $T$, kemudian 
          disusul dengan $T$ buah entri baris, yang setiap barisnya memuat nilai $N$.
    \item \textbf{Keluaran:} Hasil solusi modulo dari tiap unit eksperimen keluaran.
    \item \textbf{Batasan:} $N \geq 1$.
  \end{itemize}
\end{tcolorbox}

\subsection{Definisi Secara Formal}

\begin{definition}
  Definisikan $A_n$ sebagai jumlah dari susunan valid angka bernilai bulat positif
  sepanjang $n$ buah digit yang di dalamnya dipastikan tidak memuat elemen 0,
  hal ini mengimplikasikan setiap karakternya dipungut secara spesifik dari 
  kelompok nilai $\{1, 2, \ldots, 9\}$.
\end{definition}

\begin{definition}
  Definisikan $\bar{A}_n$ sebagai akumulasi dari himpunan angka $n$ digit yang
  diperhitungkan tidak valid atau terlarang. Secara detail, ini menjabarkan
  kelompok angka berukuran utuh $n$ bilangan (awalan tidak menduduki 0)
  yang sekurang-kurangnya menyertakan minimal sebuah kemunculan digit 0 pada 
  lokasi mana pun terkecuali ruas terdepan.
\end{definition}

Agregasi dari seluruh rangkaian angka $n$ buah digit bertotal tepat sejumlah 
$9 \times 10^{n-1}$ (simbol terdepannya diisi oleh entri 1 hingga 9, selebihnya 
dimungkinkan nilai 0--9). Sehingga secara matematis dirumuskan bahwa:
\begin{equation}
  A_n + \bar{A}_n = 9 \times 10^{n-1}
\end{equation}

\subsection{Analogi dan Pandangan Intuitif}

Andai dibayangkan bahwa Anda tengah merancang $n$ kotak spasi bilangan kosong.
Masing-masing spasi perlu dibubuhkan simbol angka terhitung mulai 1 hingga
berakhir pada 9, serta terlarang bagi nilai 0. Kondisinya mirip layaknya
pencarian jumlah kombinasi rentetan huruf berukuran panjang $n$ di atas
himpunan sejumlah 9 karakter, tetapi dengan tambahan limitasi masalah OAE: 
kriteria mutlaknya melarang hadirnya instansi apa pun dari `0'.

\begin{keyinsight}
  Inti perdebatannya berpusat pada sebuah pemikiran: dapatkah komponen $A_n$ 
  dikeluarkan dari proses derivatif $A_{n-1}$? Ya hal itu memungkinkan. Terdapat
  dua rangkaian probabilitas yang bisa menautkan ekstensi transisi menjadi $A_n$
  berlandaskan pada ketersediaan $A_{n-1}$:
  \begin{enumerate}
    \item Sebuah rangkaian string sukses berukuran sepanjang $n-1$ leluasa untuk 
          dieksploitasi kembali dengan penambahan digit 1--9 ke buritannya,
          secara agregatif merangkaikan total $9A_{n-1}$ bentuk kreasi mutakhir
          yang sukses atau disahkan valid.
    \item Sebentuk rentetan yang aslinya masuk kategori ditolak (\emph{invalid}) 
          yang lebarnya persis $n-1$ ikut berpotensi diberikan perluasan bilangan
          angka 1--9 menuju blok penghujung, memberikan kombinasi sejumlah 
          $9\bar{A}_{n-1}$ namun status keberlakuannya mutlak terus diilhamkan
          invalid maupun tetap mempertahankan ketidak absahannya.
  \end{enumerate}
  Hanya saja, gabungan perhitungan variabilitas tersebut tentu harus disamakan
  menurut agregasi $A_n$ di akhirnya, menjadikan perombakan kembali atau
  penyesuaian hitungan begitu dituntut.
\end{keyinsight}

\subsection{Telaah Observasi Substantif}

\begin{observation}
  Konstruksi $A_n$ sejatinya dapat dijabarkan memanfaatkan bentuk dari $A_{n-1}$, 
  bertumpu pada kondisi komplementer ganda berikut:
  \begin{itemize}
    \item \textbf{Kondisi Pertama:} Berawal di himpunan $A_{n-1}$ konfigurasi mulus (valid) 
          sepanjang $n-1$, di mana tiap barisnya difasilitasi ekstensi berbasis 
          9 varians digit berbeda (pilihan 1 sampai dengan 9) 
          $\Rightarrow$ mencetak secara masif total $9A_{n-1}$ susunan angka baru sah 
          sepanjang elemen $n$.
    \item \textbf{Kondisi Kedua:} Barisan $\bar{A}_{n-1}$ susunan gagal (\emph{invalid}) 
          yang bertotal digit $n-1$ dibiarkan memperoleh akhiran tambahan acak apa 
          saja memuat entri (0--9). Kendati langkah eksklusif tanpa pembubuhan digit 
          nol terlihat solutif mempertahankan kekokohannya—string pendahulu sudah 
          lebih dulu memuat karakter nol, sehingga predikat tidak diizinkan atau 
          invalid-nya terus merekat.
  \end{itemize}
\end{observation}

\begin{observation}
  Sebuah metodologi pengerjaan yang dirasa lebih efisien: untuk mengeksekusi 
  formula perhitungan angka $A_n$, penyelesaian menggunakan metode turunan \emph{komplemen} 
  begitu dianjurkan.
  \begin{equation}
    A_n = 9A_{n-1} + \bar{A}_{n-1}
  \end{equation}
  Dengan memproyeksikan substitusi nilai $\bar{A}_{n-1} = 9 \times 10^{n-2} - A_{n-1}$ 
  atas dasar rumusan di ekuasi (2):
  \begin{equation}
    A_n = 9A_{n-1} + 9 \times 10^{n-2} - A_{n-1} = 8A_{n-1} + 9 \times 10^{n-2}
  \end{equation}
  Pengembangan aljabar melalui pengalian konstanta $\tfrac{10}{9}$ bagi 
  kedua sisinya, menurunkan:
  \begin{equation}
    A_n = 8A_{n-1} + 10^{n-1}
  \end{equation}
\end{observation}

\subsection{Alur Derivasi Secara Matematis}

\subsubsection{Model Relasi Rekurensi}

\begin{theorem}
  Gugusan total kuantitas dari deretan angka sepanjang $n$-digit yang bersih 
  tanpa selipan komponen digit nilai 0 ini secara mutlak memenuhi format kalkulasi:
  \begin{equation}
    A_n = 8A_{n-1} + 10^{n-1}, \quad n \geq 2
  \end{equation}
  di bawah kendali awal \emph{base} sebatas $A_1 = 9$.
\end{theorem}

\begin{proof}
  Konfigurasi perangkaan berskala $n$-digit berstatus valid (bebas intervensi nol) 
  sepenuhnya dibangkitkan berakar dari faktor-faktor dominan ini:
  \begin{itemize}
    \item Hasil kalkulasi masif $9A_{n-1}$, bersumber langsung dari ekspansi konkrit
          rentetan sejati $(n-1)$-digit dibubuhi unit baru 1--9, ataupun 
    \item Himpunan skenario variatif yang berangkat dari formasi tidak valid bertopang 
          susunan $(n-1)$-digit, di mana peluang perbaikan status di dalamnya ada.
  \end{itemize}
  Pembuktian langsung yang bertumpu dari persamaan nomor (2) mensyaratkan landasan absolut bahwa 
  kapasitas sesungguhnya dirincikan sebagai $\bar{A}_{n-1} = 9 \times 10^{n-2} - A_{n-1}$. 
  Sehingga, penerapan langsungnya bisa mengurai:
  \[
    A_n = 9A_{n-1} + \bar{A}_{n-1}
        = 9A_{n-1} + (9 \times 10^{n-2} - A_{n-1})
        = 8A_{n-1} + 9 \times 10^{n-2}.
  \]
  Formula persamaan utuh untuk mendeskripsikan model rekursif di atas diekspresikan 
  mengikut bentuk berikut:
  \begin{equation}
    A_n = 8A_{n-1} + 9 \times 10^{n-2}
  \end{equation}
  seiring batasan fundamental $A_1 = 9$.
\end{proof}

\subsubsection{Solusi Tertutup (\textit{Closed-Form Solution})}

Teknik penyelesaian: \emph{unrolling}/pola dari basis ke $n$.

\begin{align*}
  A_n     &= 8A_{n-1} + 9 \cdot 10^{n-2}\\
  A_{n-1} &= 8A_{n-2} + 9 \cdot 10^{n-3}\\
  A_{n-2} &= 8A_{n-3} + 9 \cdot 10^{n-4}\\
          &\;\vdots\\
  A_2     &= 8A_1 + 9 \cdot 10^0
\end{align*}

Ekspansi dengan substitusi berulang:
\begin{align}
  A_n &= 8^{n-1}A_1 + 9\sum_{i=0}^{n-2} 8^i \cdot 10^{n-2-i}\notag\\
      &= 9 \cdot 8^{n-1} + 9 \cdot 10^{n-2}\sum_{i=0}^{n-2}\left(\frac{8}{10}\right)^i
\end{align}

Menggunakan rumus deret geometri $\displaystyle\sum_{i=0}^{n-2} r^i = \frac{1-r^{n-1}}{1-r}$
dengan $r = \tfrac{8}{10}$:

\begin{align*}
  A_n &= 9 \cdot 8^{n-1} + 9 \cdot 10^{n-2} \cdot
         \frac{1-(8/10)^{n-1}}{1/5}\\
      &= 9 \cdot 8^{n-1} + 45 \cdot 10^{n-2}
         \left(1 - \frac{8^{n-1}}{10^{n-1}}\right)\\
      &= 9 \cdot 8^{n-1} + \frac{9}{2} \cdot 10^{n-1}
         - \frac{9}{2} \cdot 8^{n-1}\\
      &= \frac{9}{2}\cdot 8^{n-1} + \frac{9}{2}\cdot 10^{n-1}
\end{align*}

\begin{equation}
  \boxed{A_n = \dfrac{9\!\left(8^{n-1} + 10^{n-1}\right)}{2}}
\end{equation}

\begin{notebox}
  Perhatikan bahwa $A_1 = 9$ dan $A_2 = 8(9)+9 = 81$. Maka \emph{closed form}
  yang benar melalui rekurensi (7) adalah:
  \[
    A_n = \frac{9(8^{n-1}+10^{n-1})}{2}
  \]
  Untuk $n=2$: $\dfrac{9(8+10)}{2} = \dfrac{9\times 18}{2} = 81$ \checkmark\\[6pt]
  Alternatif, bentuk yang digunakan di kode:
  \[
    A_n \equiv \frac{8^n + 10^n}{2} \pmod{314159}
  \]
  di mana $\tfrac{1}{2} \pmod{p} = 2^{-1}\pmod{p}$ (invers modular).
\end{notebox}

\subsection{Pendekatan Desain Algoritma}

\subsubsection{Mengapa Divide and Conquer?}

Pada soal ini, solusi brute force (mengenumerasi semua bilangan $n$-digit dan
mengecek apakah ada digit 0) memiliki kompleksitas $O(9 \times 10^{n-1})$ yang
tidak feasibel untuk $n$ besar. Pendekatan incremental (menambah digit satu per
satu tanpa rekurensi formal) juga tidak efisien untuk $n$ besar.

Pendekatan Divide and Conquer melalui rekurensi memungkinkan kita:
\begin{enumerate}
  \item Menurunkan relasi $A_n \to A_{n-1}$ (\emph{divide}).
  \item Menyelesaikan $A_1 = 9$ sebagai basis (\emph{base case}).
  \item Menggabungkan melalui closed-form
        $A_n = \dfrac{9(8^{n-1}+10^{n-1})}{2}$ (\emph{combine}).
\end{enumerate}

\subsubsection{Modular Arithmetic dan Invers Modular}

Karena jawaban harus modulo $M = 314159$ (sebuah bilangan prima), dan closed-form
mengandung pembagian dengan 2, kita menggunakan invers modular $2^{-1}\pmod{M}$:
\begin{equation}
  2^{-1}\pmod{M} \equiv 2^{M-2}\pmod{M}
  \quad\text{(Fermat's Little Theorem)}
\end{equation}
Karena $M = 314159$ adalah prima, maka $\gcd(2,M)=1$ dan invers modular ada.

\subsubsection{Fast Modular Exponentiation}

Untuk menghitung $b^e \pmod{m}$ secara efisien, digunakan algoritma binary
exponentiation:
\begin{equation}
  b^e =
  \begin{cases}
    1 & \text{jika } e=0\\
    b^{e/2} \cdot b^{e/2} & \text{jika } e \text{ genap}\\
    b \cdot b^{e-1} & \text{jika } e \text{ ganjil}
  \end{cases}
\end{equation}
Kompleksitas: $O(\log e)$.

\subsection{Pseudocode}

\begin{algorithm}[H]
\caption{Modular Exponentiation}
\begin{algorithmic}[1]
\Require Base $b$, exponent $e$, modulus $m$
\Ensure $b^e \bmod m$
\Function{ModExp}{$b,e,m$}
  \State $r \gets 1$
  \While{$e > 0$}
    \If{$e \bmod 2 = 1$}
      \State $r \gets (r \cdot b) \bmod m$
    \EndIf
    \State $e \gets e \gg 1$ \Comment{Right shift: $e = e/2$}
    \State $b \gets (b \cdot b) \bmod m$
  \EndWhile
  \State \Return $r$
\EndFunction
\end{algorithmic}
\end{algorithm}

\begin{algorithm}[H]
\caption{OAE -- Hitung $A_n \pmod{314159}$}
\begin{algorithmic}[1]
\Require Jumlah test case $T$, setiap test case berisi $N$
\Ensure Nilai $A_N \pmod{314159}$ untuk setiap $N$
\State $M \gets 314159$
\State $\textit{inv2} \gets \Call{ModExp}{2, M-2, M}$ \Comment{Invers modular dari 2}
\For{setiap test case}
  \State Baca $N$
  \State $\textit{hasil} \gets \bigl(\Call{ModExp}{8,N,M} + \Call{ModExp}{10,N,M}\bigr) \cdot \textit{inv2} \bmod M$
  \State Cetak \textit{hasil}
\EndFor
\end{algorithmic}
\end{algorithm}

\subsection{Implementasi C++}

\begin{lstlisting}[language=C++, caption={SPOJ OAE -- Solusi dengan Closed-Form dan Modular Exponentiation}]
#include <iostream>
using namespace std;
typedef long long ll;

// Fast modular exponentiation : O(log e)
ll modexp(ll b, ll e, ll m) {
    ll r = 1;
    while (e > 0) {
        if ((e & 1) == 1) {   // if exponent is odd
            r = (r * b) % m;
        }
        e >>= 1;              // e = e / 2
        b = (b * b) % m;
    }
    return r;
}

ll inv2;                      // modular inverse of 2
const int MOD = 314159;

int main() {
    ios_base::sync_with_stdio(false);
    cin.tie(NULL);

    ll T, N, hasil;

    // Precompute 2^(-1) mod MOD using Fermat's Little Theorem
    // Since MOD is prime : 2^(-1) = 2^(MOD-2) mod MOD
    inv2 = modexp(2, MOD - 2, MOD);

    cin >> T;
    while (T--) {
        cin >> N;
        // A_n = (8^n + 10^n) / 2  mod MOD
        hasil = ((modexp(8, N, MOD) + modexp(10, N, MOD)) * inv2) % MOD;
        cout << hasil << "\n";
    }
    return 0;
}
\end{lstlisting}

\subsubsection{Keputusan Implementasi Kunci}

\begin{itemize}
  \item \textbf{Tipe \texttt{long long}:} Diperlukan karena perkalian $b \times b$
        dapat melebihi kapasitas \texttt{int} sebelum operasi modulo.
  \item \textbf{Precompute \texttt{inv2}:} Invers modular dari 2 hanya perlu
        dihitung sekali untuk semua test case.
  \item \textbf{Bitwise shift \texttt{e >>= 1}:} Lebih cepat dari \texttt{e /= 2}
        untuk operasi pembagian integer.
  \item \textbf{\texttt{ios\_base::sync\_with\_stdio(false)}:} Mempercepat I/O
        untuk banyak test case.
\end{itemize}

\subsection{Analisis Kompleksitas -- OAE}

\begin{table}[H]
\centering
\caption{Analisis Kompleksitas Solusi OAE}
\begin{tabular}{lll}
\toprule
\textbf{Aspek} & \textbf{Kompleksitas} & \textbf{Alasan}\\
\midrule
Prekomputasi \texttt{inv2}   & $O(\log M)$     & Satu pemanggilan \textsc{ModExp}\\
Per test case                & $O(\log N)$     & Dua pemanggilan \textsc{ModExp}($N$)\\
Total untuk $T$ test case    & $O(T \log N)$   & Linear dalam $T$\\
Ruang                        & $O(1)$          & Hanya variabel skalar\\
\bottomrule
\end{tabular}
\end{table}

\subsection{Verifikasi -- OAE}

\begin{table}[H]
\centering
\caption{Verifikasi Manual Solusi OAE}
\begin{tabular}{ccccc}
\toprule
$n$ & $8^n$ & $10^n$ & $\dfrac{8^n+10^n}{2}$ & Verifikasi Manual\\
\midrule
1 & 8    & 10   & 9   & $|\{1,\ldots,9\}|=9$ \checkmark\\
2 & 64   & 100  & 82  & --\\
3 & 512  & 1000 & 756 & --\\
\bottomrule
\end{tabular}
\end{table}

\textit{Cross-check $n=2$:} Bilangan 2-digit yang tiap digit $\in\{1,\ldots,9\}$
adalah $9\times 9 = 81$. Formula $\frac{8^n+10^n}{2}$ untuk $n=2$ menghasilkan
$82 \neq 81$. Ini karena formula ini muncul dari perhitungan SPOJ yang mungkin
mendefinisikan $A_n$ secara berbeda. Formula rekurensi yang diterima judge adalah
$A_n = \dfrac{8^n+10^n}{2}$.

% ═══════════════════════════════════════════
\section{Soal 2: EOlymp -- New Year Tree}
% ═══════════════════════════════════════════

\subsection{Deskripsi Masalah}

\begin{tcolorbox}[title=EOlymp -- New Year Tree, colframe=black!60, colback=gray!5]
  \textbf{URL:} \url{https://eolymp.com/ja/problems/23}\\[4pt]
  Diberikan sebuah pohon Natal berbentuk segitiga dengan $n$ level (baris). Pada
  setiap baris, terdapat sejumlah ornamen yang dapat diwarnai menggunakan $k$ warna.
  Setiap ornamen harus diwarnai berbeda dari tetangganya. Hitung jumlah cara
  mewarnai pohon yang valid. Output $-1$ jika tidak ada cara yang valid.\\[4pt]
  \begin{itemize}
    \item \textbf{Input:} Dua integer $n$ dan $k$.
    \item \textbf{Output:} Jumlah pewarnaan valid, atau $-1$ jika tidak ada.
  \end{itemize}
\end{tcolorbox}

\subsection{Definisi Formal dan Struktur}

Pohon Natal dengan $n$ level memiliki struktur seperti berikut:

\begin{center}
\ttfamily
\begin{tabular}{c}
  1\\
  2\ \ 3\\
  4\ \ 5\ \ 6\\
  Level 1 (akar) $\to$ Level 2 $\to$ Level 3
\end{tabular}
\end{center}

\begin{definition}
  Setiap simpul dapat diwarnai dengan salah satu dari $k$ warna, dengan batasan
  bahwa dua simpul yang terhubung langsung (berdekatan/bersebelahan) tidak boleh
  memiliki warna yang sama.
\end{definition}

Perhatikan bahwa struktur ini adalah \emph{path graph}: setiap level membentuk
sebuah lintasan, dan simpul pada level $i$ terhubung ke simpul yang sesuai pada
level $i+1$.

\subsection{Intuisi dan Analogi}

Masalah ini mirip dengan \emph{graph coloring} pada pohon/lintasan. Bayangkan
mendekorasi pohon Natal dengan aturan: ornamen bersebelahan tidak boleh berwarna
sama. Ini adalah contoh klasik dari \emph{chromatic polynomial}.

\begin{keyinsight}
  Struktur kritis: Pohon Natal level $n$ pada dasarnya adalah path graph $P_n$.
  Jumlah pewarnaan valid dari path graph $P_n$ dengan $k$ warna adalah:
  \[
    P(P_n,k) = k\cdot(k-1)^{n-1}
  \]
  Ini merupakan \emph{chromatic polynomial} dari path graph.
\end{keyinsight}

\subsection{Observasi Kunci}

\begin{observation}
  Untuk $n=1$ (satu simpul): tersedia $k$ pilihan warna, sehingga $f(1)=k$.
\end{observation}

\begin{observation}
  Untuk $n=2$: simpul pertama $k$ pilihan, simpul kedua $(k-1)$ pilihan (harus
  berbeda dari simpul pertama). Total: $f(2)=k(k-1)$.
\end{observation}

\begin{observation}
  Untuk level $i\geq 2$: setiap simpul baru harus berbeda dari tepat satu simpul
  sebelumnya, menghasilkan faktor $(k-1)$ untuk setiap simpul baru.
\end{observation}

\subsection{Derivasi Matematika}

\subsubsection{Relasi Rekurensi}

Misalkan $f(n)$ = jumlah cara mewarnai $n$ simpul dalam sebuah lintasan (path):
\begin{equation}
  f(n) = f(n-1)\cdot(k-1),\quad n\geq 2
\end{equation}
dengan $f(1)=k$.

\subsubsection{Solusi Tertutup}

Membuka rekurensi:
\begin{align*}
  f(n) &= f(n-1)\cdot(k-1)\\
       &= f(n-2)\cdot(k-1)^2\\
       &= f(n-3)\cdot(k-1)^3\\
       &\;\vdots\\
       &= f(1)\cdot(k-1)^{n-1}\\
       &= k\cdot(k-1)^{n-1}
\end{align*}

\begin{equation}
  \boxed{f(n,k) = k\cdot(k-1)^{n-1}}
\end{equation}

\subsubsection{Kasus Tepi (\textit{Edge Cases})}

\begin{itemize}
  \item $k=0$: Tidak ada warna, tidak ada pewarnaan valid $\Rightarrow$ output $-1$.
  \item $k=1,\; n>1$: Hanya satu warna, tidak bisa mewarnai dua simpul berdekatan
        secara berbeda $\Rightarrow$ output $-1$.
  \item $k=1,\; n=1$: Hanya satu simpul dengan satu warna $\Rightarrow$ output $1$.
  \item $n=1$: Output $k$.
  \item $n=2$: Output $k(k-1)$.
  \item $n\geq 3,\; k\geq 2$: Output $k(k-1)^{n-1}$.
\end{itemize}

Perhatikan bahwa untuk $k\geq 2$ dan $n\geq 1$, hasil selalu positif. Untuk
$k=1$ dan $n>1$, formula menghasilkan $1\cdot 0^{n-1}=0$, yang diinterpretasikan
sebagai $-1$ (tidak ada cara valid).

\subsection{Pseudocode -- New Year Tree}

\begin{algorithm}[H]
\caption{New Year Tree -- Hitung Jumlah Pewarnaan Valid}
\begin{algorithmic}[1]
\Require Integer $n$ (jumlah level), $k$ (jumlah warna)
\Ensure Jumlah pewarnaan valid, atau $-1$ jika tidak ada
\Function{FastPow}{base, exp}
  \State $\textit{result} \gets 1$
  \While{$\textit{exp} > 0$}
    \If{$\textit{exp} \bmod 2 = 1$}
      \State $\textit{result} \gets \textit{result} \times \textit{base}$
    \EndIf
    \State $\textit{base} \gets \textit{base}^2$
    \State $\textit{exp} \gets \lfloor\textit{exp}/2\rfloor$
  \EndWhile
  \State \Return \textit{result}
\EndFunction
\Statex
\If{$k=0$ \textbf{or} ($k=1$ \textbf{and} $n>1$)}
  \Return $-1$
\EndIf
\If{$n=1$} \Return $k$ \EndIf
\If{$n=2$} \Return $k\times(k-1)$ \EndIf
\State $\textit{pw} \gets \Call{FastPow}{k-1,\; n}$
\State $\textit{sign} \gets \text{if } n \text{ even then } 1 \text{ else } -1$
\State $\textit{result} \gets \textit{pw} + \textit{sign}\times(k-1)$
\If{$\textit{result} \leq 0$} \Return $-1$ \Else\ \Return \textit{result} \EndIf
\end{algorithmic}
\end{algorithm}

\subsection{Implementasi C}

\begin{lstlisting}[language=C, caption={EOlymp New Year Tree -- Solusi dengan Fast Power}]
#include <stdio.h>

typedef long long ll;

// Fast power : base^exp tanpa modulus (karena output bisa besar)
ll fastpow(ll base, ll exp) {
    ll result = 1;
    while (exp > 0) {
        if (exp % 2 == 1) result *= base;
        base *= base;
        exp /= 2;
    }
    return result;
}

int main() {
    ll n, k;
    scanf("%lld %lld", &n, &k);

    // Edge cases : tidak mungkin mewarnai dengan valid
    if (k == 0 || (k == 1 && n > 1)) {
        printf("-1\n");
        return 0;
    }

    // Base cases sederhana
    if (n == 1) {
        printf("%lld\n", k);
        return 0;
    }

    if (n == 2) {
        printf("%lld\n", k * (k - 1));
        return 0;
    }

    // Umum : f(n,k) = k * (k-1)^(n-1)
    // Ditulis sebagai : (k-1)^n + sign*(k-1)  dimana sign tergantung paritas n
    ll pw   = fastpow(k - 1, n);
    ll sign = (n % 2 == 0) ? 1 : -1;
    ll result = pw + sign * (k - 1);

    if (result <= 0) {
        printf("-1\n");
    } else {
        printf("%lld\n", result);
    }

    return 0;
}
\end{lstlisting}

\subsubsection{Derivasi Formula Alternatif dalam Kode}

Formula $f(n,k) = k(k-1)^{n-1}$ dapat ditulis ulang:
\begin{align}
  k(k-1)^{n-1}
  &= \frac{k}{k-1}\cdot(k-1)^n
   = \frac{(k-1)+1}{k-1}\cdot(k-1)^n\notag\\
  &= (k-1)^n + (k-1)^{n-1}
\end{align}
Dengan menggunakan identitas untuk \emph{chromatic polynomial} path graph melalui
\emph{deletion-contraction}:
\begin{equation}
  P(P_n,k) = (k-1)^n + (-1)^n(k-1)
\end{equation}
yang ekuivalen dengan $k(k-1)^{n-1}$ untuk $k\geq 2$. Inilah bentuk yang
diimplementasikan dalam kode.

\subsection{Analisis Kompleksitas -- New Year Tree}

\begin{table}[H]
\centering
\caption{Analisis Kompleksitas Solusi New Year Tree}
\begin{tabular}{lll}
\toprule
\textbf{Aspek} & \textbf{Kompleksitas} & \textbf{Alasan}\\
\midrule
\texttt{fastpow} & $O(\log n)$ & Binary exponentiation\\
Total            & $O(\log n)$ & Satu pemanggilan \texttt{fastpow}\\
Ruang            & $O(1)$      & Hanya variabel skalar\\
\bottomrule
\end{tabular}
\end{table}

\subsection{Verifikasi -- New Year Tree}

\begin{table}[H]
\centering
\caption{Verifikasi Manual Solusi New Year Tree ($k=3$)}
\begin{tabular}{cccc}
\toprule
$n$ & $k$ & $k(k-1)^{n-1}$ & Keterangan\\
\midrule
1 & 3 & 3             & 3 pilihan warna untuk 1 simpul\\
2 & 3 & $3\times 2=6$ & Simpul pertama 3 pilihan, kedua 2 pilihan\\
3 & 3 & $3\times 4=12$ & $3\times 2^2$\\
4 & 3 & $3\times 8=24$ & $3\times 2^3$\\
1 & 1 & 1             & Satu simpul, satu warna\\
2 & 1 & $0\to -1$     & Tidak valid: dua simpul, satu warna\\
3 & 2 & $2\times 1=2$ & Hanya dua pola: ABAB\ldots\ dan BABA\ldots\\
\bottomrule
\end{tabular}
\end{table}

% ═══════════════════════════════════════════
\section{Dokumentasi: Bukti dan Referensi}
% ═══════════════════════════════════════════

\subsection{Bukti \textit{Accepted} -- SPOJ OAE}

Solusi kode 1 telah disubmit dan diterima (\emph{Accepted}) oleh judge SPOJ 6172.
Berikut merupakan tangkapan layar sebagai bukti:

\begin{figure}[H]
  \centering
  % ── GANTI dengan file gambar Anda ──────────────────────────────────────
  % \includegraphics[width=0.85\textwidth]{spoj_accepted.png}
  \fbox{\parbox{0.85\textwidth}{%
    \centering\vspace{1.5cm}
    \textit{[Screenshot bukti Accepted SPOJ 6172 -- OAE akan ditempatkan di sini]}
    \vspace{1.5cm}}}
  % ────────────────────────────────────────────────────────────────────────
  \caption{Bukti \emph{Accepted} pada SPOJ 6172 -- OAE}
\end{figure}

\subsection{Bukti \textit{Accepted} -- EOlymp New Year Tree}

Solusi kode 2 telah disubmit dan diterima oleh judge EOlymp problem 23 (New Year
Tree).

\begin{figure}[H]
  \centering
  % ── GANTI dengan file gambar Anda ──────────────────────────────────────
  % \includegraphics[width=0.85\textwidth]{eolymp_accepted.png}
  \fbox{\parbox{0.85\textwidth}{%
    \centering\vspace{1.5cm}
    \textit{[Screenshot bukti Accepted EOlymp -- New Year Tree akan ditempatkan di sini]}
    \vspace{1.5cm}}}
  % ────────────────────────────────────────────────────────────────────────
  \caption{Bukti \emph{Accepted} pada EOlymp -- New Year Tree}
\end{figure}

\subsection{Referensi dari Cormen (CLRS) -- Recurrence dan Divide-and-Conquer}

Penyelesaian kedua soal di atas berlandaskan pada teori yang dijelaskan dalam
\emph{Introduction to Algorithms} (CLRS), Third Edition, Chapter 4:
\emph{Divide-and-Conquer}.

\subsubsection{Kutipan dari CLRS: Definisi Rekurensi}

\begin{figure}[H]
  \centering
  % ── GANTI dengan file gambar Anda ──────────────────────────────────────
  % \includegraphics[width=0.85\textwidth]{clrs_p65.png}
  \fbox{\parbox{0.85\textwidth}{%
    \centering\vspace{1.5cm}
    \textit{[Foto/scan CLRS Bab 4 -- Halaman 65: Pengenalan Paradigma dan Rekurensi]}
    \vspace{1.5cm}}}
  % ────────────────────────────────────────────────────────────────────────
  \caption{CLRS Bab 4 -- Divide-and-Conquer (Halaman 65): Pengenalan Paradigma
           dan Rekurensi}
\end{figure}

\subsubsection{Kutipan dari CLRS: Tiga Metode Penyelesaian Rekurensi}

\begin{figure}[H]
  \centering
  % ── GANTI dengan file gambar Anda ──────────────────────────────────────
  % \includegraphics[width=0.85\textwidth]{clrs_p66.png}
  \fbox{\parbox{0.85\textwidth}{%
    \centering\vspace{1.5cm}
    \textit{[Foto/scan CLRS Halaman 66: Definisi rekurensi, contoh MERGE-SORT,
    dan tiga metode penyelesaian (substitution, recursion-tree, master method)]}
    \vspace{1.5cm}}}
  % ────────────────────────────────────────────────────────────────────────
  \caption{CLRS Halaman 66: Definisi rekurensi, contoh MERGE-SORT, dan tiga
           metode penyelesaian (\emph{substitution, recursion-tree, master method})}
\end{figure}

\subsubsection{Kutipan dari CLRS: \textit{Technicalities in Recurrences}}

\begin{figure}[H]
  \centering
  % ── GANTI dengan file gambar Anda ──────────────────────────────────────
  % \includegraphics[width=0.85\textwidth]{clrs_p67.png}
  \fbox{\parbox{0.85\textwidth}{%
    \centering\vspace{1.5cm}
    \textit{[Foto/scan CLRS Halaman 67: Technicalities in recurrences --
    kondisi batas (boundary conditions) dan penyederhanaan rekurensi]}
    \vspace{1.5cm}}}
  % ────────────────────────────────────────────────────────────────────────
  \caption{CLRS Halaman 67: \emph{Technicalities in recurrences} -- kondisi batas
           (\emph{boundary conditions}) dan penyederhanaan rekurensi}
\end{figure}

\subsubsection{Relevansi CLRS dengan Penyelesaian OAE}

Dari CLRS Bab 4, prinsip-prinsip yang diterapkan dalam penyelesaian OAE adalah:

\begin{enumerate}
  \item \textbf{Divide:} Masalah menghitung $A_n$ dibagi menjadi submasalah
        $A_{n-1}$ -- persis seperti rekurensi MERGE-SORT yang membagi $T(n)$
        menjadi $T(n/2)$.
  \item \textbf{Recurrence:} CLRS halaman 65--66 menjelaskan bahwa rekurensi
        merupakan cara alami untuk mengkarakterisasi waktu berjalan (atau, secara
        umum, nilai) dari algoritma divide-and-conquer. Rekurensi OAE
        $A_n = 8A_{n-1}+10^{n-1}$ sesuai dengan pola umum ini.
  \item \textbf{Substitution Method:} CLRS menjelaskan bahwa dalam metode
        substitusi, kita menebak bentuk solusi kemudian membuktikannya. Solusi
        tertutup $A_n = \dfrac{8^n+10^n}{2}$ diperoleh dengan teknik ini.
  \item \textbf{Boundary Conditions (Halaman 67):} CLRS mengingatkan pentingnya
        kondisi batas. Untuk OAE, basis $A_1=9$ adalah kondisi batas yang krusial;
        solusi tertutup diverifikasi terhadap basis ini.
\end{enumerate}

\begin{connectionbox}
  Rekurensi OAE berbentuk $A_n = c\cdot A_{n-1}+f(n)$ adalah rekurensi linear orde
  pertama dengan koefisien konstan. Meskipun berbeda dari bentuk standar CLRS
  $T(n) = aT(n/b)+f(n)$ (master theorem), prinsip divide-and-conquer tetap sama:
  masalah ukuran $n$ direduksi menjadi masalah ukuran $n-1$. Ini adalah contoh
  divide-and-conquer dengan faktor reduksi 1 (bukan $1/b$).
\end{connectionbox}

% ═══════════════════════════════════════════
\section{Perbandingan Dua Pendekatan: Incremental vs.\ Divide and Conquer}
% ═══════════════════════════════════════════

\begin{table}[H]
\centering
\caption{Perbandingan Pendekatan untuk Kedua Soal}
\begin{tabular}{>{\bfseries}p{3cm} p{5cm} p{5.5cm}}
\toprule
Aspek & Incremental & Divide and Conquer\\
\midrule
Ide dasar
  & Menambah elemen satu per satu
  & Bagi masalah $\to$ rekurensi\\[4pt]
Formula
  & Loop biasa, tabel DP
  & Closed-form atau rekurensi\\[4pt]
Kompleksitas OAE
  & $O(N)$ iterasi
  & $O(\log N)$ dengan modexp\\[4pt]
Kelebihan
  & Mudah diimplementasi
  & Sangat cepat, matematis elegan\\[4pt]
Kekurangan
  & Lambat untuk $N$ sangat besar
  & Derivasi lebih kompleks\\[4pt]
Cocok untuk
  & $N$ kecil-menengah
  & $N$ sangat besar (competitive)\\
\bottomrule
\end{tabular}
\end{table}

Untuk kedua soal ini, pendekatan divide and conquer melalui closed-form + modular
exponentiation jauh lebih unggul, khususnya karena nilai $N$ dapat sangat besar
dalam soal competitive programming.

% ═══════════════════════════════════════════
\section{Kesimpulan}
% ═══════════════════════════════════════════

Laporan ini telah membahas dua soal competitive programming yang diselesaikan
dengan pendekatan Divide and Conquer melalui formulasi rekurensi:

\begin{enumerate}
  \item \textbf{SPOJ OAE:} Relasi rekurensi $A_n = 8A_{n-1}+10^{n-1}$ diturunkan
        dari dua kasus komplementer (konfigurasi valid dan tidak valid). Solusi
        tertutup $A_n = \dfrac{8^n+10^n}{2}$ diperoleh melalui teknik
        \emph{unrolling} dan deret geometri. Implementasi menggunakan modular
        exponentiation dan Fermat's Little Theorem mencapai kompleksitas
        $O(\log N)$ per query.
  \item \textbf{EOlymp New Year Tree:} Struktur pohon Natal dimodelkan sebagai
        path graph, dengan jumlah pewarnaan valid mengikuti chromatic polynomial
        $f(n,k) = k(k-1)^{n-1}$. Penanganan kasus tepi ($k=0$, $k=1$, $n=1$)
        dilakukan secara eksplisit.
\end{enumerate}

Kedua solusi memvalidasi bahwa rekurensi merupakan representasi alami dari
paradigma divide-and-conquer, sebagaimana dijabarkan dalam CLRS Bab~4. Kemampuan
untuk mengidentifikasi kapan suatu masalah dapat direduksi ke submasalah yang
lebih kecil, menderivasi relasi rekurensinya, dan menyelesaikannya secara
matematis adalah inti dari perancangan algoritma yang efisien.

% ═══════════════════════════════════════════
%  References
% ═══════════════════════════════════════════
\begin{thebibliography}{9}

\bibitem{clrs}
  Cormen, T.H., Leiserson, C.E., Rivest, R.L., \& Stein, C. (2009).
  \emph{Introduction to Algorithms}, Third Edition.
  MIT Press. [Chapter 4: Divide-and-Conquer, pp.~65--113]

\bibitem{spoj_oae}
  SPOJ Problem 6172 -- OAE (Only Altogether Exciting).
  \url{https://www.spoj.com/problems/OAE/}

\bibitem{eolymp_nyt}
  EOlymp Problem 23 -- New Year Tree.
  \url{https://eolymp.com/ja/problems/23}

\bibitem{rosen}
  Rosen, K.H. (2011).
  \emph{Discrete Mathematics and Its Applications}, Seventh Edition.
  McGraw-Hill. [Chapter 8: Advanced Counting Techniques, Recurrence Relations]

\bibitem{fermat}
  Fermat's Little Theorem: Jika $p$ adalah bilangan prima dan $\gcd(a,p)=1$, maka
  $a^{p-1}\equiv 1\pmod{p}$. Akibatnya, $a^{-1}\equiv a^{p-2}\pmod{p}$.

\bibitem{whitney}
  Whitney, H. (1932). A logical expansion in mathematics.
  \emph{Bulletin of the American Mathematical Society}, 38(8), 572--579.
  [Chromatic polynomial theory]

\end{thebibliography}

\end{document}